\problem{P07}{缺失相关系数的范围}{相关矩阵补全；构造}
三个随机变量 $X,Y,Z$ 均已标准化。已知 $\Corr(X,Y)=a$、$\Corr(X,Z)=b$，其中 $|a|,|b|<1$。求 $c=\Corr(Y,Z)$ 的全部可行值，并为每个可行值构造联合分布。代入 $a=3/5,b=4/5$，解释区间端点的退化关系。
\solution 令 $U=Y-aX$、$V=Z-bX$。它们与 $X$ 不相关，且
\[
 \Var(U)=1-a^2,\quad\Var(V)=1-b^2,
 \quad\Cov(U,V)=c-ab.
\]
Cauchy--Schwarz 不等式给出必要条件
\[
 \boxed{|c-ab|\le\sqrt{(1-a^2)(1-b^2)}}.
\]
反过来，设 $q=(c-ab)/\sqrt{(1-a^2)(1-b^2)}\in[-1,1]$。取独立标准正态 $G_0,G_1,G_2$，并令
\[
 \begin{aligned}
 X&=G_0,\\
 Y&=aG_0+\sqrt{1-a^2}G_1,\\
 Z&=bG_0+\sqrt{1-b^2}\bigl(qG_1+\sqrt{1-q^2}G_2\bigr).
 \end{aligned}
\]
直接计算即得三个方差均为一且相关系数符合要求，故条件充分。$|q|<1$ 时构造矩阵可逆，相关矩阵正定；$q=\pm1$ 时两个残差线性相关，矩阵奇异。
数值情形给出 $ab=12/25$、根号项也是 $12/25$，故 $0\le c\le24/25$。下端点满足 $V=-3U/4$，上端点满足 $V=3U/4$；这说明达到端点不是相关矩阵失效，而是三个变量之间出现精确线性关系。
\insight 剔除共同变量以后，未知相关性变成两个残差的协方差；必要界与可达构造可以在同一坐标系中完成。
